%========================================
% MatinBook — Stage 06: Theorem System Test
% Version: v1.1
% Purpose:
%   Validate the theorem system of MatinBook v1.1:
%     1. All theorem-like environments (theorem, lemma,
%        corollary, proposition, definition, example,
%        remark, exercise, solution, proof)
%     2. Both box families (matinbox and matin-outline)
%     3. Shared counter behavior
%     4. Cross-references with \cref
%     5. Index entries inside boxes
%
% Compilation:
%   xelatex -shell-escape stage06-theorems.tex
%   xelatex -shell-escape stage06-theorems.tex
%   xelatex -shell-escape stage06-theorems.tex
%   (3 passes recommended for index)
%
% Expected outcome:
%   A professionally typeset PDF demonstrating that every
%   theorem environment renders correctly, that the shared
%   counter produces continuous numbering, and that
%   cross-references resolve to the correct names.
%========================================

%------------------------------------------------------------------------
% Search Path (for tests directory)
%------------------------------------------------------------------------
\makeatletter
\def\input@path{{../}}
\makeatother

%------------------------------------------------------------------------
% Document Class
%------------------------------------------------------------------------
\documentclass{matinbook}

%------------------------------------------------------------------------
% Book Metadata
%------------------------------------------------------------------------
\title{آزمون سیستم قضایا و محیط‌های علمی}
\author{تیم توسعه MatinBook}
\date{مهر ۱۴۰۵}

\booktitle{آزمون سیستم قضایا و محیط‌های علمی}

\renewcommand{\repository}{github.com/matinmahmoudi-cs/matinbook}

%========================================================================
% DOCUMENT
%========================================================================
\begin{document}

%------------------------------------------------------------------------
% FRONT COVER
%------------------------------------------------------------------------
\makecover
    {آزمون سیستم قضایا و محیط‌های علمی}
    {ارزیابی محیط‌های علمی، جعبه‌ها و ارجاع هوشمند}
    {تیم توسعه MatinBook}
    {مهر ۱۴۰۵}

%------------------------------------------------------------------------
% FRONT MATTER (symmetric margins)
%------------------------------------------------------------------------
\frontmatter
\frontmattergeometry

% Title page
\maketitle

% Copyright / Colophon
\thispagestyle{empty}
\vfill
\begin{center}
    {\large\textbf{شناسنامه آزمون}}

    \vspace{0.8cm}

    \begin{tabular}{rl}
        عنوان: & آزمون سیستم قضایا و محیط‌های علمی \\
        نسخه: & \matinbookversion \\
        تاریخ: & \matinbookdate \\
        موتور: & \lr{XeLaTeX} \\
        مجوز: & \lr{MIT} \\
    \end{tabular}

    \vspace{0.8cm}

    {\small این سند بخشی از مجموعه آزمون‌های یکپارچگی \lr{MatinBook} است.}
\end{center}
\clearpage

% Preface
\chapter*{پیشگفتار}
\addcontentsline{toc}{chapter}{پیشگفتار}

این سند، ششمین آزمون از مجموعه‌ی آزمون‌های یکپارچگی
\lr{MatinBook} نسخه‌ی \lr{\matinbookversion} است. هدف آن،
اعتبارسنجی کامل سیستم محیط‌های علمی کلاس در پنج حوزه‌ی
متمایز است: محیط‌های قضیه‌ای، خانواده‌های جعبه، شمارنده‌ی
مشترک، ارجاع هوشمند با \lr{\textbackslash cref}، و
ورودی‌های نمایه داخل جعبه‌ها.

در کتاب‌های ریاضی و فنی، محیط‌های علمی ستون فقرات
ساختاردهی محتوا هستند. قضیه، لم، نتیجه، گزاره، تعریف،
مثال، نکته، تمرین، راه‌حل و اثبات، هر یک نقش معنایی
مشخصی دارند و باید در قالب بصری متمایز نمایش داده شوند.
\lr{MatinBook} با دو خانواده‌ی جعبه‌ی \lr{matinbox} و
\lr{matin-outline}، این تمایز را ایجاد می‌کند.

هر بخش از این آزمون، با مثال‌های عملی و ارجاعات متقابل
همراه است تا خواننده بتواند درستی رفتار هر زیرسیستم را
در بافت واقعی یک کتاب ارزیابی کند.

\clearpage

% Table of Contents
\tableofcontents
\clearpage

% List of Figures
\listoffigures
\clearpage

% List of Tables
\listoftables
\clearpage

%------------------------------------------------------------------------
% MAIN MATTER (wide margin for notes)
%------------------------------------------------------------------------
\mainmatter
\mainmattergeometry

%========================================================================
% CHAPTER 1 — NUMBER THEORY
%========================================================================
\chapter{نظریه اعداد}
\label{chap:numbertheory}

\section{اعداد اول}
\label{sec:numbertheory-primes}

\begin{definition}[عدد اول]
    عدد طبیعی $n > 1$ را \textbf{اول}\index{عدد اول} گوییم
    هرگاه تنها مقسوم‌علیه‌های مثبت آن ۱ و خود $n$ باشند.
    اعدادی که اول نباشند را \textbf{مرکب} گوییم.
    \label{def:prime}
\end{definition}

طبق \cref{def:prime}، عدد ۱ نه اول است و نه مرکب. این
قرارداد برای سادگی در بیان قضایا اتخاذ شده است.

\mnote{عدد ۱ نه اول است و نه مرکب. این قرارداد در نظریه اعداد استاندارد است.}

\begin{example}
    اعداد اول کوچک‌تر از ۵۰ عبارت‌اند از:
    \[
    2, 3, 5, 7, 11, 13, 17, 19, 23, 29, 31, 37, 41, 43, 47
    \]
    \label{ex:primes-under-50}
\end{example}

\begin{example}
    عدد ۶ اول نیست (مرکب است) زیرا $6 = 2 \times 3$.
    مقسوم‌علیه‌های عدد ۶ عبارت‌اند از: ۱، ۲، ۳، ۶.
    \label{ex:not-prime}
\end{example}

\section{قضیه‌ی اساسی حساب}
\label{sec:numbertheory-fundamental}

\begin{theorem}[قضیه‌ی اساسی حساب]
    هر عدد طبیعی $n > 1$ را می‌توان به‌طور یکتا
    (صرف‌نظر از ترتیب عوامل) به‌صورت حاصل‌ضرب اعداد اول نوشت:
    \[
    n = p_{1}^{e_{1}} \cdot p_{2}^{e_{2}} \cdots p_{k}^{e_{k}}
    \]
    که در آن $p_{1} < p_{2} < \cdots < p_{k}$ اعداد اول
    و $e_{i} \geq 1$ اعداد صحیح هستند.
    \label{thm:fundamental}
\end{theorem}

طبق \cref{thm:fundamental}، هر عدد طبیعی بزرگ‌تر از ۱
دارای یک تجزیه‌ی یکتا به عوامل اول است. برای مثال:
\[
60 = 2^{2} \times 3 \times 5
\]

\begin{proof}
    ابتدا \textbf{وجود} تجزیه را با استقرای قوی اثبات می‌کنیم.

    \textbf{پایه:} عدد ۲ اول است و تجزیه‌ی آن خودش می‌باشد.

    \textbf{گام استقرا:} فرض کنید هر عدد $m$ با $2 \leq m < n$
    قابل تجزیه به عوامل اول باشد. برای $n$ دو حالت داریم:

    \begin{enumerate}
        \item اگر $n$ اول باشد، تجزیه‌ی آن خودش است.
        \item اگر $n$ مرکب باشد، $n = a \cdot b$ که
        $2 \leq a, b < n$. طبق فرض استقرا، $a$ و $b$
        هر یک به عوامل اول تجزیه می‌شوند. از ضرب این دو
        تجزیه، تجزیه‌ی $n$ به‌دست می‌آید.
    \end{enumerate}

    برای اثبات \textbf{یکتایی} از لم اقلیدس
    (\cref{lem:euclid}) استفاده می‌شود که در ادامه
    خواهیم دید.
\end{proof}

\section{لم اقلیدس}
\label{sec:numbertheory-euclid}

\begin{lemma}[لم اقلیدس]
    اگر $p$ عددی اول باشد و $p \mid ab$، آنگاه
    $p \mid a$ یا $p \mid b$.
    \label{lem:euclid}
\end{lemma}

\begin{proof}
    فرض کنید $p \nmid a$. آنگاه $\gcd(p, a) = 1$
    (چون $p$ اول است). طبق قضیه‌ی بزو (\lr{Bezout})،
    اعداد صحیح $x$ و $y$ وجود دارند که:
    \[
    px + ay = 1
    \]
    با ضرب طرفین در $b$:
    \[
    pbx + aby = b
    \]
    از آنجا که $p \mid p$ و $p \mid ab$، نتیجه می‌شود
    $p \mid b$.
\end{proof}

\begin{corollary}
    اگر $p$ عددی اول باشد و $p \mid a_{1}a_{2} \cdots a_{n}$،
    آنگاه $p$ حداقل یکی از $a_{i}$ها را می‌شمارد
    (بر آن بخش‌پذیر است).
    \label{cor:euclid-general}
\end{corollary}

\begin{proposition}
    کوچک‌ترین مقسوم‌علیه بزرگ‌تر از ۱ هر عدد طبیعی $n > 1$،
    یک عدد اول است.
    \label{prop:smallest-divisor}
\end{proposition}

\section{نکات مهم}
\label{sec:numbertheory-remarks}

\begin{remark}
    لم اقلیدس (\cref{lem:euclid}) یکی از پایه‌های اساسی
    نظریه‌ی اعداد است و در اثبات بسیاری از قضایا از جمله
    یکتایی تجزیه در قضیه‌ی اساسی حساب
    (\cref{thm:fundamental}) به کار می‌رود.
    \label{rem:euclid-importance}
\end{remark}

\begin{remark}
    عدد ۲ تنها عدد اول زوج است. تمام اعداد اول دیگر فرد
    هستند، زیرا اگر یک عدد زوج بزرگ‌تر از ۲ اول می‌بود،
    باید بر ۲ بخش‌پذیر می‌بود که با تعریف عدد اول در
    تناقض است.
    \label{rem:two-is-special}
\end{remark}

%========================================================================
% CHAPTER 2 — GROUP THEORY
%========================================================================
\chapter{مبانی نظریه‌ی گروه‌ها}
\label{chap:grouptheory}

\section{تعریف گروه}
\label{sec:grouptheory-definition}

\begin{definition}[گروه]
    یک \textbf{گروه}\index{گروه} $(G, *)$ شامل یک مجموعه‌ی
    ناتهی $G$ و یک عمل دوتایی $*$ روی $G$ است که در چهار
    شرط زیر صدق کند:

    \begin{enumerate}
        \item \textbf{بسته بودن:}
        $\forall a, b \in G: a * b \in G$
        \item \textbf{شرکت‌پذیری:}
        $\forall a, b, c \in G: (a * b) * c = a * (b * c)$
        \item \textbf{وجود عضو خنثی:}
        $\exists e \in G: \forall a \in G: e * a = a * e = a$
        \item \textbf{وجود عضو وارون:}
        $\forall a \in G: \exists a^{-1} \in G:
        a * a^{-1} = a^{-1} * a = e$
    \end{enumerate}
    \label{def:group}
\end{definition}

\begin{example}
    مجموعه‌ی اعداد صحیح $\Z$ با عمل جمع معمولی یک گروه است:

    \begin{itemize}
        \item عضو خنثی: $0$
        \item وارون $a$: $-a$
    \end{itemize}
    \label{ex:z-group}
\end{example}

\begin{example}
    مجموعه‌ی اعداد حقیقی غیرصفر $\R \setminus \{0\}$ با
    عمل ضرب معمولی یک گروه است:

    \begin{itemize}
        \item عضو خنثی: $1$
        \item وارون $a$: $\frac{1}{a}$
    \end{itemize}
    \label{ex:r-group}
\end{example}

\section{قضایای مقدماتی}
\label{sec:grouptheory-basic}

\begin{theorem}[یکتایی عضو خنثی]
    عضو خنثی در هر گروه یکتاست.
    \label{thm:unique-identity}
\end{theorem}

\begin{proof}
    فرض کنید $e$ و $e'$ هر دو عضو خنثی گروه $G$ باشند.
    آنگاه:
    \[
    e = e * e' \quad \text{(چون $e'$ خنثی است)}
    \]
    \[
    e * e' = e' \quad \text{(چون $e$ خنثی است)}
    \]
    بنابراین $e = e'$ و عضو خنثی یکتاست.
\end{proof}

\begin{theorem}[یکتایی عضو وارون]
    وارون هر عضو در یک گروه یکتاست.
    \label{thm:unique-inverse}
\end{theorem}

\begin{proof}
    فرض کنید $a \in G$ و $b$ و $c$ دو وارون $a$ باشند:

    \begin{align*}
        b &= b * e \quad \text{(خنثی)} \\
        &= b * (a * c) \quad \text{($c$ وارون $a$ است)} \\
        &= (b * a) * c \quad \text{(شرکت‌پذیری)} \\
        &= e * c \quad \text{($b$ وارون $a$ است)} \\
        &= c \quad \text{(خنثی)}
    \end{align*}

    بنابراین $b = c$ و عضو وارون یکتاست.
\end{proof}

\section{تمرین‌ها}
\label{sec:grouptheory-exercises}

\begin{exercise}
    ثابت کنید که در هر گروه $G$، معادله‌ی $a * x = b$
    همواره دارای جواب یکتاست.
    \label{exc:group-equation}
\end{exercise}

\begin{solution}
    \textbf{وجود جواب:} $x = a^{-1} * b$ یک جواب است زیرا:
    \[
    a * (a^{-1} * b) = (a * a^{-1}) * b = e * b = b
    \]

    \textbf{یکتایی:} فرض کنید $x_{1}$ و $x_{2}$ دو جواب باشند:
    \begin{align*}
        a * x_{1} &= b \\
        a * x_{2} &= b
    \end{align*}
    با برابر قرار دادن و ضرب از چپ در $a^{-1}$:
    \[
    a^{-1} * (a * x_{1}) = a^{-1} * (a * x_{2})
    \implies x_{1} = x_{2}
    \]
    بنابراین جواب یکتاست.
\end{solution}

\begin{exercise}
    ثابت کنید که $(\Z_{n}, +_{n})$ یک گروه است (مجموعه‌ی
    اعداد صحیح به پیمانه‌ی $n$ با عمل جمع پیمانه‌ای).
    \label{exc:zn-group}
\end{exercise}

%========================================================================
% CHAPTER 3 — CROSS-REFERENCES TEST
%========================================================================
\chapter{آزمون ارجاع متقابل}
\label{chap:cref}

\section{ارجاع به تمام محیط‌ها}
\label{sec:cref-all}

در این بخش، توانایی ارجاع به انواع مختلف محیط‌های علمی
را آزمایش می‌کنیم:

\begin{itemize}
    \item \textbf{قضیه:} \cref{thm:fundamental} می‌گوید هر
    عدد طبیعی بزرگ‌تر از ۱ به‌صورت حاصل‌ضرب اعداد اول
    قابل تجزیه است.

    \item \textbf{تعریف:} طبق \cref{def:prime}، عدد اول
    عددی است که دقیقاً دو مقسوم‌علیه مثبت داشته باشد.

    \item \textbf{لم:} \cref{lem:euclid} یک لم کلیدی در
    نظریه‌ی اعداد است.

    \item \textbf{نتیجه:} \cref{cor:euclid-general} تعمیم
    لم اقلیدس به چند عامل است.

    \item \textbf{گزاره:} \cref{prop:smallest-divisor}
    کوچک‌ترین مقسوم‌علیه را معرفی می‌کند.

    \item \textbf{مثال:} \cref{ex:primes-under-50} اعداد
    اول کوچک‌تر از ۵۰ را فهرست کرده است.

    \item \textbf{نکته:} \cref{rem:euclid-importance}
    اهمیت لم اقلیدس را توضیح می‌دهد.

    \item \textbf{تمرین:} \cref{exc:group-equation} حل
    معادله در گروه را می‌طلبد.
\end{itemize}

\section{ارجاع بین فصلی}
\label{sec:cref-cross-chapter}

در \cref{chap:analysis} (فصل بعدی)، مفاهیم حد و پیوستگی
را بررسی خواهیم کرد. همچنین در \cref{exc:zn-group}
(فصل قبل)، از شما خواسته شد گروه $\Z_{n}$ را بررسی کنید.

\section{ارجاع گروهی و بازه‌ای}
\label{sec:cref-group}

می‌توان به چند مرجع به‌طور هم‌زمان ارجاع داد:
\cref{thm:unique-identity,thm:unique-inverse} هر دو به
یکتایی در گروه اشاره دارند.

%========================================================================
% CHAPTER 4 — REAL ANALYSIS
%========================================================================
\chapter{مبانی آنالیز حقیقی}
\label{chap:analysis}

\section{حد و همگرایی}
\label{sec:analysis-limit}

\begin{definition}[حد دنباله]
    دنباله‌ی $(a_{n})_{n=1}^{\infty}$ به عدد حقیقی $L$
    \textbf{همگرا} است اگر:
    \[
    \forall \varepsilon > 0, \;
    \exists N \in \N, \;
    \forall n \geq N: \abs{a_{n} - L} < \varepsilon
    \]
    در این صورت می‌نویسیم $\lim_{n \to \infty} a_{n} = L$.
    \label{def:limit}
\end{definition}

\begin{example}
    دنباله‌ی $a_{n} = \frac{1}{n}$ به صفر همگراست، زیرا:
    \[
    \forall \varepsilon > 0, \;
    \text{با انتخاب } N > \frac{1}{\varepsilon}, \;
    \forall n \geq N: \abs{\frac{1}{n} - 0} < \varepsilon
    \]
    \label{ex:convergent-1over}
\end{example}

\begin{theorem}[یکتایی حد]
    اگر دنباله‌ای همگرا باشد، حد آن یکتاست.
    \label{thm:unique-limit}
\end{theorem}

\begin{proof}
    فرض کنید $\lim a_{n} = L_{1}$ و $\lim a_{n} = L_{2}$
    با $L_{1} \neq L_{2}$. قرار دهید
    $\varepsilon = \frac{\abs{L_{1} - L_{2}}}{2} > 0$.

    طبق تعریف حد (\cref{def:limit})، عدد $N_{1}$ وجود دارد
    که برای $n \geq N_{1}$، $\abs{a_{n} - L_{1}} < \varepsilon$.
    همچنین عدد $N_{2}$ وجود دارد که برای $n \geq N_{2}$،
    $\abs{a_{n} - L_{2}} < \varepsilon$.

    برای $n \geq \max(N_{1}, N_{2})$:
    \begin{align*}
        \abs{L_{1} - L_{2}}
        &= \abs{(L_{1} - a_{n}) + (a_{n} - L_{2})} \\
        &\leq \abs{L_{1} - a_{n}} + \abs{a_{n} - L_{2}} \\
        &< \varepsilon + \varepsilon = \abs{L_{1} - L_{2}}
    \end{align*}
    که تناقض است. بنابراین $L_{1} = L_{2}$.
\end{proof}

\begin{remark}
    عکس \cref{thm:unique-limit} لزوماً برقرار نیست: یکتایی
    حد شرط کافی برای همگرایی نیست. دنباله‌ی
    $a_{n} = (-1)^{n}$ حد ندارد، با اینکه اگر داشت
    یکتا می‌بود.
    \label{rem:limit-converse}
\end{remark}

%========================================================================
% CHAPTER 5 — SUMMARY
%========================================================================
\chapter{خلاصه‌ی نتایج}
\label{chap:summary}

\section{آمار محیط‌های استفاده شده}
\label{sec:summary-environments}

جدول \cref{tab:environments-used} محیط‌های علمی استفاده‌شده
در این آزمون را نشان می‌دهد.

\begin{matintable}{محیط‌های علمی استفاده‌شده در این آزمون}{tab:environments-used}
\begin{tblr}{
    width = \textwidth,
    colspec = {l X[c] X[c]},
    hlines, vlines,
    colsep = 8pt,
    rowsep = 4pt,
    row{1} = {font=\bfseries},
}
محیط & تعداد & وضعیت \\
\lr{theorem} & ۴ & \textcolor{matingreen}{فعال} \\
\lr{definition} & ۳ & \textcolor{matingreen}{فعال} \\
\lr{lemma} & ۱ & \textcolor{matingreen}{فعال} \\
\lr{corollary} & ۱ & \textcolor{matingreen}{فعال} \\
\lr{proposition} & ۱ & \textcolor{matingreen}{فعال} \\
\lr{example} & ۵ & \textcolor{matingreen}{فعال} \\
\lr{remark} & ۳ & \textcolor{matingreen}{فعال} \\
\lr{proof} & ۵ & \textcolor{matingreen}{فعال} \\
\lr{exercise} & ۲ & \textcolor{matingreen}{فعال} \\
\lr{solution} & ۱ & \textcolor{matingreen}{فعال} \\
\end{tblr}
\end{matintable}

\section{معیارهای ارزیابی}
\label{sec:summary-criteria}

در این آزمون، سیستم محیط‌های علمی \lr{MatinBook} بر اساس
پنج معیار ارزیابی شد:

\begin{enumerate}
    \item \textbf{ظاهر بصری}: تمایز دو خانواده‌ی جعبه.
    \item \textbf{شمارنده‌ی مشترک}: شماره‌گذاری پیوسته.
    \item \textbf{ارجاع هوشمند}: \lr{\textbackslash cref} با نام‌های فارسی.
    \item \textbf{شکست صفحات}: رفتار جعبه‌های بلند.
    \item \textbf{نمایه}: ورودی‌های \lr{\textbackslash index} داخل جعبه‌ها.
\end{enumerate}

\section{نتیجه‌گیری}
\label{sec:summary-conclusion}

\textbf{تمام محیط‌های علمی \lr{MatinBook} نسخه‌ی
\lr{\matinbookversion} با موفقیت بارگذاری و آزمایش شدند.
ده نوع جعبه‌ی رنگی، شمارنده‌ی مشترک، ارجاع هوشمند و
اثبات‌های ریاضی، همگی به‌درستی کار می‌کنند.}

%------------------------------------------------------------------------
% BACK MATTER (symmetric margins)
%------------------------------------------------------------------------
\backmatter
\frontmattergeometry

% Bibliography (empty placeholder)
\printbibliography[title={منابع و مراجع}]

% Index
\printindex

%------------------------------------------------------------------------
% BACK COVER
%------------------------------------------------------------------------
\authorbio{%
    تیم توسعه‌ی \lr{MatinBook}، گروهی از پژوهشگران و برنامه‌نویسان
    فارسی‌زبان است که با هدف ارائه‌ی یک راه‌حل حرفه‌ای برای
    حروف‌چینی کتاب‌های فنی فارسی فعالیت می‌کند. این پروژه حاصل
    سال‌ها تجربه در حوزه‌ی نشر دیجیتال و برنامه‌نویسی است.
}

\makebackcover{%
    این سند، ششمین آزمون از مجموعه‌ی آزمون‌های یکپارچگی
    \lr{MatinBook} نسخه‌ی \lr{\matinbookversion} است.

    \vspace{0.5cm}

    \textbf{زیرسیستم‌های آزموده‌شده:}
    \begin{itemize}
        \item محیط‌های قضیه‌ای (۱۰ نوع)
        \item دو خانواده‌ی جعبه
        \item شمارنده‌ی مشترک
        \item ارجاع هوشمند
        \item ورودی‌های نمایه
    \end{itemize}
}

\end{document}